\section{Theoretical analysis: what closure does and does not buy}\label{sec:theory}
Write $\mathfrak a=\spn(A)$, $G=B^\top B$, $P_0$ for the orthogonal projector onto $\mathfrak a$ and $N=I-P_0$.
All statements below are proved in Appendix~\ref{app:proofs} and checked numerically by unit tests in the
supplementary code.

\paragraph{Invariances of the objective.}
With $\varepsilon=0$, a change of generator coordinates $B\mapsto BM$, $M\in GL(K)$, gives
$\Lc(BM)=K^{-2}\|R(M\otimes M)\|_F^2$, where $R$ is the residual matrix. Hence
$\sigma_{\min}(M)^4\Lc\le\Lc(BM)\le\sigma_{\max}(M)^4\Lc$.

The loss is therefore invariant exactly under orthogonal $M$, it is not invariant under general invertible
$M$, and its zero set, ``$\mathfrak a$ is a subalgebra'', is $GL(K)$-invariant. For $\varepsilon>0$ the
regularized projector satisfies $P_\varepsilon(BM)=B(G+\varepsilon M^{-\top}M^{-1})^{-1}B^\top$, so it is
not even a function of the span. A genuinely span-only objective evaluates $\Lc$ on the orthonormal basis
$\sqrt2\,BG^{-1/2}$; we test it as a separate variant (Sec.~\ref{sec:ablations}).

Because each generator is normalized to $\sqrt2$, $\operatorname{diag}G=2$. The Gram penalty therefore acts
only on off-diagonal entries, and its target ($2I$ or $I$) does not affect gradients.

\begin{proposition}[Exact closure integrates]\label{prop:integrate}
If $A_1,\dots,A_K\in\so(d)$ are linearly independent and $[A_i,A_j]\in\mathfrak a$ for all $i,j$, then
$\mathfrak a$ is a $K$-dimensional Lie subalgebra. It is the Lie algebra of a unique connected, immersed (not
necessarily closed) subgroup $H\subseteq SO(d)$ generated by $\exp(\mathfrak a)$.
\end{proposition}

Let $\delta(\mathfrak a)=\sup_{X,Y\in\mathfrak a\setminus 0}\|N[X,Y]\|_F/(\|X\|_F\|Y\|_F)$ be the uniform
escape constant.

\begin{lemma}[The loss controls uniform escape]\label{lem:delta}
$\delta(\mathfrak a)\le K\sqrt{\Lc(A;0)}\,/\,\lambda_{\min}(G)$. For an orthogonal basis normalized to
$\sqrt2$, this is $\tfrac K2\sqrt{\Lc}$.
\end{lemma}

\begin{proposition}[Approximate closure controls BCH escape]\label{prop:bch}
For $X,Y\in\mathfrak a$ with $\|X\|_F,\|Y\|_F\le r<\tfrac{\ln 2}{4}$,
\[
\dist\!\big(\log(e^Xe^Y),\mathfrak a\big)\;\le\;\tfrac{\delta}{4}\big[-\ln(2-e^{4r})-4r\big]
\;\le\;\tfrac12\delta r^2+C(r_0)\,\delta r^3\quad(r\le r_0),
\]
where $C(r_0)=\big(\tfrac14[-\ln(2-e^{4r_0})-4r_0]-4r_0^2\big)/r_0^3$.
\end{proposition}

Both bounds vanish under exact closure, consistent with Proposition~\ref{prop:integrate}. The radius is
small: it covers rotations by at most about $0.12$\,rad. The bound is therefore a local guarantee about the
learned family near the identity, not about every action the coefficient network can output.

\begin{proposition}[Chained orthogonal stability]\label{prop:chain}
If $z_{t+1}=R_tz_t+e_t$ and $\hat z_{t+1}=R_t\hat z_t$ with $R_t\in SO(d)$, then
$\|\hat z_T-z_T\|\le\|\hat z_0-z_0\|+\sum_{t<T}\|e_t\|$.
\end{proposition}

\paragraph{Closure is not recovery.}
We separate four properties, each implied by the next:
(i)~closure;
(ii)~the correct conjugacy class of subalgebra of $\so(d)$;
(iii)~ground-truth recovery in the learned coordinates, i.e.\ $\spn(R\mathfrak aR^\top)=\mathfrak g$ after the
latent alignment $R$ determined by the data;
(iv)~additionally, correct action inference.

\begin{proposition}[Closed subalgebras in the benchmarks]\label{prop:classify}
(a)~Every non-degenerate 2-dimensional subalgebra of $\so(5)$ is abelian and conjugate to the true
$\mathfrak t^2$. Hence, for $T^2$, (i)$\Rightarrow$(ii).
(b)~Every 3-dimensional subalgebra of $\so(4)\cong\su(2)\oplus\su(2)$ is conjugate either to the diagonal
$\so(3)$ acting on $\mathbb R^3\oplus\mathbb R$ (the true action) or to a \emph{chiral} factor $\su(2)_{L}$ or
$\su(2)_{R}$. These factors are exchanged by a reflection and are not conjugate to the diagonal $\so(3)$.
Hence, for SO(3), (i)$\not\Rightarrow$(ii).
\end{proposition}

The two SO(3) classes have identical abstract structure but differ in representation. The conjugation-invariant
participation ratio $\operatorname{tr}(X^\top X)^2/\operatorname{tr}((X^\top X)^2)$ of unit elements is $2$
for the diagonal class and $4$ for a chiral factor. The mean commutator norm differs as well ($\sqrt2$ vs.\
$2$ under our normalization), so a commutator threshold identifies only abelian versus non-abelian algebras,
not the representation. No step of the objective forces (iii) or (iv). We claim no identifiability result.
